\documentclass[12pt,a4paper]{article}
\usepackage[T2A,T1]{fontenc}
\usepackage[utf8]{inputenc}
\usepackage[english,russian]{babel}
\usepackage{XCharter}
\usepackage{amsmath,amssymb,mathtools}
\usepackage{array,longtable,booktabs}
\usepackage[margin=24mm]{geometry}
\usepackage[hidelinks]{hyperref}
\urlstyle{same}
\setlength{\parindent}{0pt}
\setlength{\parskip}{0.45em}
\setlength{\LTleft}{0pt}
\setlength{\LTright}{0pt}
\emergencystretch=2em
\setlength{\tabcolsep}{4pt}
\newcolumntype{P}[1]{>{\raggedright\arraybackslash}p{#1}}
\title{Корпус работ Н. П. Брусенцова: что доступно и как читать}
\author{}
\date{}
\begin{document}
\maketitle

\emph{Исследовательская заметка. Проверка доступности и скачивание: 20.11.2026.}

\section{Что удалось найти}

Основной каталог полных текстов размещён на \href{https://ternarycomp.cs.msu.ru/all.html}{сайте троичной информатики МГУ}; выходные данные сверялись с \href{https://ternarycomp.cs.msu.ru/Brusentsov/Brusentsov_v2.html}{библиографией Н. П. Брусенцова на сайте МГУ}. Дополнительный ранний материал 1971 года и архивный файл статьи 2003 года найдены на сайте \href{https://computer-museum.ru/histussr/trilog0.htm}{Виртуального компьютерного музея}.

В репозиторий добавлены 18 PDF-файлов и исходный Word-документ статьи 2003 года. Большинство файлов скачано непосредственно с сайта МГУ. Для статьи 2003 года сохранён оригинальный ZIP-архив музея и извлечённый из него DOC; PDF рядом — удобная копия, пересобранная LibreOffice. Она не заменяет опубликованный экземпляр.

\textbf{Полная книга «Реставрация логики» не найдена в открытом каталоге.} Скачана отдельная глава Н. П. Брусенцова «Реанимация аристотелевой силлогистики» (2005, с. 140–145). Не следует ссылаться на неё как на полный текст книги.

\section{Как устроено предложение Н. П. Брусенцова}

Чтобы не смешивать разные логические объекты, сначала фиксируем, что именно обозначено в его формулах.

\begin{itemize}
\item Малые буквы \(x\), \(y\) обозначают термины/признаки; формулы также рассматривают классы вещей, которым они присущи. Термин, класс его носителей и целое высказывание \(P\) — разные уровни обозначения.
\item Сочетание \(xy\) обозначает сочетание признаков, \(xy'\) — \(x\) с противоположностью \(y\), \(x'y\) и \(x'y'\) — два оставшихся сочетания. Штрих у термина обозначает противоположный признак; конкретный класс предметов в каждой записи зависит от заданной области.
\item Отношение \(x\Rightarrow y\) он предлагает понимать как необходимое следование между терминами. Его четырёхпозиционная схема \(+-0+\), в порядке \(xy, xy', x'y, x'y'\), задаёт статусы сочетаний: необходимое включение, необходимое исключение, привходящий статус, необходимое включение. Ноль здесь — не третье значение целого высказывания \(P\to Q\), а статус отношения к одной комбинации терминов. В главе 2005 года автор явно связывает нормальную форму с непустотой \(xy\) и \(x'y'\), пустотой \(xy'\) и привходящим статусом \(x'y\); профиль \(+-0+\) кодирует эти статусы в принятом универсуме. Сам знак \(+\) не является квантором, а чтение профиля на конкретном конечном списке — наша предметная иллюстрация авторской семантики. \href{../sources/pdfs/Брусенцов_Реанимация-аристотелевой-силлогистики_2005-глава-из-Реставрации-логики.pdf}{Глава 2005 года, PDF с. 1–3}.
\item В статье 2010 года запись \(x\to y\) обозначает материальную импликацию внутри алгебры терминов Н. П. Брусенцова; она не является автоматически пропозициональной формулой \(P\to Q\). Вектор \(+-++\), полученный для обычной материальной импликации по четырём предметным комбинациям, — наше сравнение, а не введённая автором векторная таблица.
\item Схема действует в принятом Н. П. Брусенцовым универсуме: для каждого первичного термина должны сосуществовать вещи, которым он присущ, и вещи, которым он противоположен. Это содержательное предположение о рассматриваемых терминах, а не следствие одной только материальной импликации.
\end{itemize}

\textbf{Тезис автора, его оговорка и граница результата.} Н. П. Брусенцов не оставляет открытым вопрос о том, тождественна ли обычная двухзначная материальная импликация необходимому содержательному следованию: он прямо отрицает такую тождественность без дополнительных условий и предлагает собственное отношение \(x\Rightarrow y\) и его статусное представление. Содержательность он объясняет тем, что качество \(y\), присущее \(x\)-вещам, необходимо содержится в их сущности; это авторское философское основание, а не независимо установленный причинный или физический факт. В четырёхпозиционной схеме \(+-0+\) ячейка \(x'y\) отмечена как привходящая. В главе 2005 года автор прямо выводит непустоту \(xy\) и \(x'y'\), пустоту \(xy'\) и привходящий статус \(x'y\) при сосуществовании противоположностей; профиль кодирует эти статусы, а его чтение на конечном списке — наша предметная иллюстрация. В статье 2010 года автор оговаривает равнозначность материальной импликации и необходимого следования в специально принятом универсуме \(УА\), подчинённом постулату сосуществования противоположностей. Аргумент главы 2005 года о строгой импликации Льюиса \(V'_{xy'}\) без условий существования \(x\) и \(y'\) следует отличать от тезиса о материальной импликации. Результат модели условен относительно определений и постулатов автора; он не доказывает общей несостоятельности классической логики или исторической тождественности модели логике Аристотеля.

\textbf{Сопоставление употребляемых кодов.} В записи \(A(x,y)\in\{1,0,\sigma\}\) знак \(\in\) означает «принадлежит множеству», а скобки перечисляют возможные значения. В статье 2003 года: \(1\) — необходимая присущность, \(0\) — антиприсущность, \(\sigma\) — привходящий статус. Условие \(0<\sigma<1\) (сигма больше нуля и меньше единицы) задаёт его положение между крайними значениями, но не фиксирует конкретное число и не уточняет, допускаются ли разные числовые реализации. В главе 2005 года обобщённая запись \(xy\lor\sigma x'y\lor x'y'\) помечает ячейку \(x'y\) символом \(\sigma\); знак \(\lor\) здесь соединяет выражения об классах, а не целые высказывания. В знаковой шкале главы привходящий статус записан нулём. Поэтому профиль \(+-0+\), в порядке \(xy,xy',x'y,x'y'\), не использует тот же смысл нуля, что \(A(x,y)\) 2003 года. Слова «необходимое включение», «необходимое исключение» и «привходящий статус» относятся к шкале модели, а не автоматически к физической или метафизической необходимости.

Профиль \(+-++\) — наше сравнение: он кодирует истинность точечной материальной импликации \(X(t)\to Y(t)\), которая ложна только для сочетания \(xy'\). Здесь \(X(t)\) означает, что отдельный предмет \(t\) имеет признак \(x\), а \(Y(t)\) — признак \(y\). В профиле \(+-++\) знак \(+\) означает истинность, знак \(-\) — ложность этой формулы; в авторском \(+-0+\) знаки обозначают статусы сочетаний классов. Совпадение \(+\) или \(-\) не означает совпадения типа статуса.

В формулах карты \(Axy\) в контексте статьи 2010 года — общеутвердительное суждение «Все \(x\) суть \(y\)»; \(A(x,y)\) в статье 2003 года — обозначение статусной функции. Их связь нужно устанавливать по контексту публикации. \(P,Q\) — целые высказывания; \(\models\) обозначает семантическое следование, \(\vdash\) — выводимость по правилам. Ни один из этих знаков не является обозначением \(x\Rightarrow y\), отношения между терминами.

Именно это различие следует проверять в статьях: Н. П. Брусенцов предлагает модель \textbf{связи между терминами/классами вещей}. Из этого ещё не следует, что он опроверг обычную импликацию как функцию от истинностных значений двух высказываний или заменил ею все современные понятия семантического следования.

\section{Основные работы о силлогистике и следовании}

\begin{longtable}{P{0.260\textwidth}P{0.300\textwidth}P{0.360\textwidth}}
\toprule
\textbf{Работа и выходные данные} & \textbf{Копия в репозитории} & \textbf{Для чего читать} \\
\midrule
\endfirsthead
\toprule
\textbf{Работа и выходные данные} & \textbf{Копия в репозитории} & \textbf{Для чего читать} \\
\midrule
\endhead
\href{https://www.computer-museum.ru/books/setun/brusencov_03.pdf}{«Математическая теория силлогистики»} (1971), \emph{Вычислительная техника и вопросы кибернетики}, вып. 8, с. 154–176 & \href{../sources/pdfs/Брусенцов_Математическая-теория-силлогистики_1971.pdf}{PDF, 19 страниц} & Ранняя математическая постановка вопроса: как выразить строение категорических суждений и правила силлогистики средствами алгебры. У каталожного диапазона с. 154–176 и числа страниц PDF есть расхождение; полноту именно этого файла ещё нужно сверить с выпуском сборника. \\
\href{https://ternarycomp.cs.msu.ru/Papers/Brusentsov_Polnaya_sistema_V19.pdf}{«Полная система категорических силлогизмов Аристотеля»} (1982), \emph{Вычислительная техника и вопросы кибернетики}, вып. 19, с. 3–17 & \href{../sources/pdfs/Брусенцов_Полная-система-категорических-силлогизмов-Аристотеля_1982.pdf}{PDF, 15 страниц} & Сверить раннюю формализацию полного набора категорических силлогизмов с более поздней трёхзначной реконструкцией. PDF — скан; страницы визуально читаемы. \\
\href{https://computer-museum.ru/histussr/trilog0.htm}{«Трехзначная интерпретация силлогистики Аристотеля»} (2003), \emph{Историко-математические исследования}, 2-я серия, вып. 8 (43), с. 317–327 & \href{../sources/legacy-docs/TRILOG33.DOC}{архивный DOC}, \href{../sources/archives/Трехзначная-интерпретация-силлогистики-Аристотеля_archive.zip}{исходный ZIP}, \href{../sources/pdfs/Брусенцов_Трехзначная-интерпретация-силлогистики-Аристотеля_2003_версия-музея.pdf}{PDF-копия}, 10 страниц после конвертации & Самый развёрнутый из скачанных текстов о модели. Автор задаёт трёхзначные статусы присущности, вводит постулат сосуществования противоположностей и описывает алгебраическое получение заключений. Конвертированная пагинация отличается от 11 страниц печатного диапазона; тождество архивной версии опубликованному тексту требует сверки. \\
\href{https://ternarycomp.cs.msu.ru/Papers/Rea.pdf}{«Реанимация аристотелевой силлогистики»} (2005), глава книги \emph{Реставрация логики}, с. 140–145 & \href{../sources/pdfs/Брусенцов_Реанимация-аристотелевой-силлогистики_2005-глава-из-Реставрации-логики.pdf}{PDF, 7 страниц} & Самый прямой короткий источник по \(+-0+\): Н. П. Брусенцов начинает со строгой импликации Льюиса, затем добавляет условия существования терминов и объясняет, почему ячейка \(x'y\) для его отношения несущественна. Это глава, а не книга целиком. \\
\href{https://ternarycomp.cs.msu.ru/Papers/Soder_sled.pdf}{«Содержательное следование, сосуществование противоположностей и формальные импликации»} (2006), материалы IX конференции «Современная логика», с. 342–343 & \href{../sources/pdfs/Брусенцов_Содержательное-следование-и-формальные-импликации_2006.pdf}{PDF, 2 страницы} & Кратко формулирует тезис о различии материальной импликации и предложенного содержательного следования. Полезен как декларация позиции автора, но аргументация в нём очень сжата. \\
\href{https://ternarycomp.cs.msu.ru/Papers/Alrek.pdf}{«Алгебраическая реконструкция силлогистики»} (2008), \emph{Информационные технологии}, № 2, с. 7–10 & \href{../sources/pdfs/Брусенцов_Алгебраическая-реконструкция-силлогистики_2008.pdf}{PDF, 7 страниц} & Даёт формулы и список типов суждений; утверждает, что в предлагаемой реконструкции получается 128 правильных модусов. Печатный диапазон и число электронных страниц различаются. \\
\href{https://ternarycomp.cs.msu.ru/Papers/Sborn9.pdf}{«О содержательном истолковании логико-алгебраических выражений»} (2008), \emph{Программные системы и инструменты}, № 9, с. 127–128 & \href{../sources/pdfs/Брусенцов_О-содержательном-истолковании-логико-алгебраических-выражений_2008.pdf}{PDF, 3 страницы} & Особенно полезен для различения объектов: в этой трактовке логические выражения описывают вещи и их признаки; пропуск комбинации может означать несущественность, а не её исключение. \\
\href{https://ternarycomp.cs.msu.ru/Papers/Arger.pdf}{«Аристотелева силлогистика и гераклитово сосуществование противоположностей»} (2009), \emph{Историко-математические исследования}, 2-я серия, вып. 13 (48), с. 270–273 & \href{../sources/pdfs/Брусенцов_Аристотелева-силлогистика-и-сосуществование-противоположностей_2009.pdf}{PDF, 4 страницы} & Объясняет принятую автором онтологическую предпосылку — сосуществование признака и его противоположности. Историческое отождествление этой предпосылки с позицией Аристотеля нужно проверять отдельно. \\
\href{https://ternarycomp.cs.msu.ru/Papers/brusentsov_sborn2010_1.pdf}{«Трехзначность отношения следования»} (2010), \emph{Программные системы и инструменты}, № 11, с. 86–87 & \href{../sources/pdfs/Брусенцов_Трехзначность-отношения-следования_2010.pdf}{PDF, 2 страницы} & Сжатое изложение трёх статусов: необходимо, возможно, невозможно. Важен для проверки, как автор связывает их с \(A\), \(I\), \(E\) и отношениями терминов. \\
\href{https://ternarycomp.cs.msu.ru/Papers/brusentsov_sborn2010_2.pdf}{«Воссоздание аристотелевой безукоризненной силлогистики»} (2010), там же, с. 88–91 & \href{../sources/pdfs/Брусенцов_Воссоздание-аристотелевой-силлогистики_2010.pdf}{PDF, 4 страницы} & Поздняя формулировка тезиса и постулата сосуществования; приводит собственное чтение текста «Первой аналитики». Сопоставлять с переводом Аристотеля, не принимать исторические выводы статьи как независимое подтверждение. \\
\href{https://ternarycomp.cs.msu.ru/Papers/brusentsov_algebraiz_2011.pdf}{«Алгебраизация безупречной логики Аристотеля путем усовершенствования кэрролова метода индексов»} (2011), фрагмент книги Н. П. Брусенцова и Ю. С. Владимировой, с. 3–5 & \href{../sources/pdfs/Брусенцов_Владимирова_Алгебраизация-силлогистики-методом-Кэрролла_2011.pdf}{PDF-фрагмент, 3 страницы} & Показывает представление следования через несовместимость \(xy'\) и элиминацию среднего термина в Barbara. Это фрагмент книги, не её полный текст. \\
\href{https://ternarycomp.cs.msu.ru/Papers/brusentsov_obobsch_alg_2011.pdf}{«Трехзначное обобщение алгебры логики. Преодоление несовершенности ДНФ…»} (2011), фрагмент той же книги, с. 5–6 & \href{../sources/pdfs/Брусенцов_Трехзначное-обобщение-булевой-алгебры_2011.pdf}{PDF-фрагмент, 2 страницы} & Самое ясное краткое различение исключённой ячейки \(xy'\) и несущественной ячейки \(x'y\). На последней странице автор также заявляет, что Аристотель ошибся в одном из разбираемых им положений — существенная оговорка к тезису о «воссоздании» его логики. \\
\href{https://ternarycomp.cs.msu.ru/Papers/brusentsov_usoversh_log_umozakl_book.pdf}{«Усовершенствование логики умозаключений»} (2012), брошюра, в библиографии 8 страниц & \href{../sources/pdfs/Брусенцов_Усовершенствование-логики-умозаключений_2012.pdf}{PDF, 6 страниц} & Позднее краткое изложение: необходимое следование задаётся повторными включениями признаков; материальная импликация отличается добавлением несущественной ячейки. PDF занимает 6 страниц при 8 страницах в библиографии; полноту проверить не удалось. \\
«Аристотелево необходимое следование» (2014), \emph{Историко-математические исследования}, 2-я серия, вып. 15 (50), с. 239–241; «Трехзначное обобщение алгебры логики…» (2014), там же, с. 241–242 (диапазоны подтверждаются \href{https://www.mathnet.ru/php/person.phtml?option_lang=rus\&personid=68115}{библиографической записью Math-Net}). \href{https://ternarycomp.cs.msu.ru/Papers/Brusen_IMI14_Neobh_sled.pdf}{Общий PDF} & \href{../sources/pdfs/Брусенцов_Аристотелево-необходимое-следование-и-ДНФ_2014.pdf}{PDF, 4 страницы} & В локальной копии первая статья видна на печатных с. 239–241; вторая начинается на с. 241, но её продолжение на с. 242 отсутствует: PDF-с. 3 обрывается внутри текста, а PDF-с. 4 содержит только публикационную помету. Копия второй статьи неполная. \\
\bottomrule
\end{longtable}

\section{Вычислительная линия}

Эти публикации показывают, как авторы кодируют и обрабатывают отношения в программной форме. Они не доказывают сами по себе, что выбранное истолкование исторически принадлежит Аристотелю или должно заменить современную логику.

\begin{longtable}{P{0.260\textwidth}P{0.300\textwidth}P{0.360\textwidth}}
\toprule
\textbf{Работа} & \textbf{Копия} & \textbf{Что исследовать} \\
\midrule
\endfirsthead
\toprule
\textbf{Работа} & \textbf{Копия} & \textbf{Что исследовать} \\
\midrule
\endhead
Н. П. Брусенцов, Ю. С. Владимирова, Х. Рамиль Альварес, \href{https://ternarycomp.cs.msu.ru/Papers/Sborn6.pdf}{«Троичный логико-алгебраический и арифметический процессор»} (2005), \emph{Программные системы и инструменты}, № 6, с. 184–187 & \href{../sources/pdfs/Брусенцов_Владимирова_Рамиль-Альварес_Троичный-логико-алгебраический-процессор_2005.pdf}{PDF, 4 страницы} & Контекст связи троичного процессора с алгебраическим аппаратом и аргументацией авторов о троичной информатике. \\
Н. П. Брусенцов и Ю. С. Владимирова, \href{https://ternarycomp.cs.msu.ru/Papers/MMRO12.pdf}{«Троичная компьютеризация логики»} (2005), 12-я конференция ММРО, с. 40–42 & \href{../sources/pdfs/Брусенцов_Владимирова_Троичная-компьютеризация-логики_2005_ММРО-12.pdf}{PDF, 2 страницы} & Утверждает, что трёхзначные шкалы кодируют функции двухзначных терминов, а третье значение появляется на уровне отношений. \\
Н. П. Брусенцов и Ю. С. Владимирова, \href{https://ternarycomp.cs.msu.ru/Papers/Sborn8.pdf}{«Конструктная компьютеризация аристотелевой силлогистики»} (2007), \emph{Программные системы и инструменты}, № 8, с. 24–28 & \href{../sources/pdfs/Брусенцов_Владимирова_Конструктная-компьютеризация-силлогистики_2007_ПСИ-8.pdf}{PDF, 5 страниц} & Версия с описанием конструктного кодирования и компьютерного вывода. \\
Н. П. Брусенцов и Ю. С. Владимирова, \href{https://ternarycomp.cs.msu.ru/Papers/MMRO13.pdf}{«Конструктная компьютеризация силлогистики»} (2007), ММРО-13, с. 10–13 & \href{../sources/pdfs/Брусенцов_Владимирова_Конструктная-компьютеризация-силлогистики_2007_ММРО-13.pdf}{PDF, 3 страницы} & Показывает операцию склеивания кодов и приводит Barbara; заявляет 128 правильных модусов. Сопоставить с числом 192 в статье 2003 года. \\
\bottomrule
\end{longtable}

\section{Что уже видно и что надо проверить}

\begin{enumerate}
\item \textbf{Это не только формализация правила вывода.} Н. П. Брусенцов задаёт онтологию объектов: признаки, классы вещей, их комбинации, существование и несуществование. Связь \(x\Rightarrow y\) определяется внутри этой модели, а не таблицей истинности для двух целых высказываний.
\item \textbf{У модели есть дополнительное условие.} Сосуществование \(x\), \(x'\), \(y\), \(y'\) исключает вакуумные случаи и позволяет из запрета \(xy'\) выводить нужные включения. Без этого условия отсутствие контрпримеров может быть вызвано пустым классом и не свидетельствовать о содержательной связи.
\item \textbf{Исторический статус — отдельный вопрос.} Н. П. Брусенцов соотносит модель с Аристотелем, но в статье 2003 года говорит, что постулат сосуществования не был явно высказан Аристотелем; в более поздних текстах он же поправляет Аристотеля в понимании некоторых отношений. Корректнее пока говорить о реконструкции Н. П. Брусенцова, вдохновлённой его чтением Аристотеля, а не о доказанно точном изложении Аристотеля.
\item \textbf{Есть внутреннее расхождение по числу модусов.} Статья 2003 года сообщает 192 правильных модуса, а статья ММРО-13 (2007) и «Алгебраическая реконструкция силлогистики» (2008) — 128. Статья ММРО-13 ссылается на публикацию 2003 года; в библиографии статьи 2008 года ссылки на ММРО-13 не обнаружено. Основания двух подсчётов в просмотренных текстах не сведены, поэтому требуется восстановить классификации и сопоставить критерии счёта.
\item \textbf{Нужна проверка против первичных и современных источников.} Следующий содержательный шаг — сверить цитаты \emph{Первая аналитика} 57b и 24b26 с греческим текстом и надёжным переводом, затем сравнить формальную модель с современными семантиками, где отдельно задаются пустые термины, модальность и релевантность. Компьютеризация доказывает реализуемость заданных операций, но не снимает эти семантические и исторические вопросы.
\end{enumerate}
\section{Дополнительные доступные материалы}

\begin{itemize}
\item \textbf{«Начала информатики» (1994, 176 страниц).} Полный текст доступен для просмотра на стороннем \href{https://djvu.online/file/mkIqfvqkmGScQ}{DjVu Online}; в каталоге МГУ книга также указана. В корпус не копировалась: это не официальный репозиторий издания.
\item \textbf{«Искусство достоверного рассуждения» (1998).} Книга важна как расширенное изложение реконструкции; библиографически подтверждена, но отдельный полный открытый текст в просмотренных каталогах не обнаружен.
\item \textbf{«Аристотелева силлогистика в символической логике Льюиса Кэрролла» (2011).} Полный текст книги в каталогах не найден; в корпусе сохранены два отдельных фрагмента (с. 3–6).
\item \textbf{«Диаграммы Льюиса Кэрролла и аристотелева силлогистика» (1977).} Библиографическая запись подтверждена; отдельно скачать полный текст не удалось. Этот пробел можно закрыть через экземпляр сборника или библиотеку.
\item \textbf{«Реставрация логики» (2005, 164 страницы).} В репозитории только глава «Реанимация аристотелевой силлогистики»; наличие остальных глав в книге не означает, что они доступны отдельными файлами.
\end{itemize}
\section{Рекомендуемый порядок чтения}

\begin{enumerate}
\item 1971 и 1982 — исходная алгебраизация и ранняя система силлогизмов.
\item 2003 — развернутая модель \(+\), \(-\), \(0\) и постулат сосуществования.
\item 2005 и 2006 — прямое различение строгой/материальной импликации и предлагаемого следования.
\item 2008–2010 — формулы, терминные отношения и силлогистические следствия.
\item 2011–2014 — фрагменты метода Кэрролла и поздняя компактная формулировка.
\item Публикации о компьютеризации — после того, как станет ясно, какие именно объекты и операции кодируются.
\end{enumerate}
Каждое утверждение об Аристотеле в корпусе следует маркировать как цитату, интерпретацию Н. П. Брусенцова или собственную проверку исследователя. Не переносить слово «истинность» между этими уровнями без оговорки: иногда автор обсуждает истинность суждения, иногда — статус комбинации терминов, а иногда — необходимость/возможность отношения.

\section{Основные источники карты}

\begin{itemize}
\item МГУ. \href{https://ternarycomp.cs.msu.ru/Brusentsov/Brusentsov_v2.html}{Список научных трудов и изобретений Н. П. Брусенцова}.
\item МГУ. \href{https://ternarycomp.cs.msu.ru/all.html}{Каталог размещённых полных текстов}.
\item Н. П. Брусенцов. \href{https://computer-museum.ru/histussr/trilog0.htm}{«Трехзначная интерпретация силлогистики Аристотеля» (2003)}.
\item Н. П. Брусенцов. \href{https://ternarycomp.cs.msu.ru/Papers/Rea.pdf}{«Реанимация аристотелевой силлогистики» (2005)}.
\item Н. П. Брусенцов. \href{https://ternarycomp.cs.msu.ru/Papers/Soder_sled.pdf}{«Содержательное следование, сосуществование противоположностей и формальные импликации» (2006)}.
\item Н. П. Брусенцов. \href{https://ternarycomp.cs.msu.ru/Papers/Alrek.pdf}{«Алгебраическая реконструкция силлогистики» (2008)}.
\item Н. П. Брусенцов, Ю. С. Владимирова. \href{https://ternarycomp.cs.msu.ru/Papers/MMRO13.pdf}{«Конструктная компьютеризация силлогистики» (ММРО-13, 2007)}.
\item \href{../sources/README.md}{Описание локального корпуса и происхождения файлов}.
\end{itemize}
\end{document}
